%============================================================
\chapter{手続き上の記録による実証：\texorpdfstring{$M(t)$}{M(t)} と改修の率}
\label{ch:protocol-emp}
%============================================================

第\ref{ch:protocol}章の区分「内」では、本稿が他の章で個社データを要するとして
測定を断念した量が個別に公開されている（第\ref{sec:protocol-observability}節）。
本章はこれを観測する。

母集団は基準時点で稼働していた $\Phi$ の全数であり、標本は事後に選んでいない。
第\ref{part:catalog}部の演繹に対する初めての観測である。

\section{検証する含意の登録}
\label{sec:protocol-predictions}

第\ref{ch:next}節の手続きに従い、取得の前に予測を固定する。
本章の予測は四つあるため、$\Phi$ に番号を付して区別する。

\begin{quote}
\textbf{予測$\Phi$1}：手続きが呼び出す他の手続きの数が少ないほど、
改修されるまでの期間が長い。命題\ref{prop:unmanned-period}の $\lambda_{\rm dep}$ 側の含意である。

\textbf{対抗仮説}：改修の頻度は依存の数ではなく利用量で決まる。
依存が少なくとも利用が多ければ改修される。
\end{quote}

\begin{quote}
\textbf{予測$\Phi$2}：取引が途絶えた $\Phi$ について、
$M(t)$ は途絶に先行して単調に減少する。
第\ref{sec:objective}節の $\tau=\inf\{t:M(t)<0\}$ が吸収状態であることの含意である。

\textbf{対抗仮説}：途絶は $M(t)$ と無関係に生じる。
$M(t)$ は途絶の直前まで変化しない。
\end{quote}

\begin{quote}
\textbf{予測$\Phi$3}：区分「内」で稼働する $\Phi$ の $\CCC$ の分布は非負に偏る。
系\ref{cor:growth-needs-capital}より、$\CCC<0$ の実現には $E>0$ を要するためである。

\textbf{対抗仮説}：$\CCC$ の分布は第\ref{part:industry}章の企業部門と同様であり、
区分「内」に固有の偏りはない。
\end{quote}

\begin{quote}
\textbf{予測$\Phi$4}：稼働する $\Phi$ を第\ref{sec:type-assignment}節の規則で割り当てると、
族7は 7-4 を除いてほぼ現れない。命題\ref{prop:beneficiary}の含意である。

\textbf{対抗仮説}：族7型も他の族と同程度の比重で現れる。
\end{quote}

\subsection{事前に確定した事項。}

\begin{center}
\begin{tabularx}{\textwidth}{@{}lX@{}}
\toprule
項目 & 内容 \\
\midrule
母集団 & 基準時点に手続きの記録上で稼働していた $\Phi$ の全数。事後に選ばない \\
下限 & 取引件数が基準期間に一定数未満のものを除く。閾値は取得前に確定する \\
途絶の定義 & 一定期間にわたり取引が生じないこと。区分「内」では手続き自体は停止しないため、停止ではなく途絶で定義する \\
定常状態 & $\CCC$ は定常を前提とするため、$r$ が安定した期間に限る（第\ref{sec:age}節） \\
観測単位 & 識別子。主体単位への集計は行わない（注\ref{rem:additivity}） \\
類型の割り当て & 第\ref{sec:type-assignment}節の判定順序を用いる。順序は変更しない \\
\bottomrule
\end{tabularx}
\captionof{table}{予測$\Phi$1--$\Phi$4 について事前に確定した事項。}
\label{tab:protocol-registration}
\end{center}

\begin{remark}[$\lambda_{\rm dep}$ と利用量は相関しうる]
\label{rem:dep-usage-confound}
予測$\Phi$1 の対抗仮説は棄却しにくい。
依存の多い手続きは機能が豊富であり、利用量も大きい傾向があるためである。
検定では利用量を統制する必要があるが、
統制後も残差に依存の効果が現れるかは事前には分からない。
\textbf{統制の手続きを取得前に確定しておく。}
\end{remark}

\section{検定の結果}
\label{sec:protocol-results}

公開型の分散台帳の記録により四つを検定した。
母集団は基準時点で稼働していた $\Phi$ の全数であり、標本は事後に選んでいない。
支持されたのは一つである。

\begin{center}
\begin{tabularx}{\textwidth}{@{}llXl@{}}
\toprule
予測 & 標本 & 結果 & 判定 \\
\midrule
$\Phi$1 & 改修可能な手続き 193 万件 & 依存の数は改修までの期間を説明しない & 棄却 \\
$\Phi$2 & 手数料と残高の双方を持つ 427 件 & 途絶前 30 日で減少（$p=3.2\times10^{-9}$） & 支持 \\
$\Phi$3 & 稼働する $\Phi$ 52 件 & 対抗仮説が識別されない & 検証不能 \\
$\Phi$4 & 同上 & 族7 は 1 件だが 7-2 であり 7-4 は現れない & 棄却 \\
\bottomrule
\end{tabularx}
\captionof{table}{予測$\Phi$1--$\Phi$4 の検定結果。}
\label{tab:protocol-results}
\end{center}

\begin{remark}[$\Phi$1：複製の集中が推定を支配する]
\label{rem:clone-concentration}
193 万の識別子が指す実装は 76,609 種しかない。
うち一つを 155 万件、すなわち母集団の 80\% が指している。
これらは独立な $\Phi$ ではなく、一つの $\Phi$ の複製である。

全体では依存の数の係数が正で有意になるが（$p=0.004$）、
この一つのクラスタを外すと符号が逆転して有意でなくなる（$p=0.45$）。
\textbf{クラスタ頑健標準誤差では防げない。}
クラスタ内の相関は補正するが、一つのクラスタが曝露の 95\% を占める場合の
点推定の偏りは補正しないためである。

注\ref{rem:additivity}の加法性の喪失は主体と識別子の乖離を述べたが、
同一の実装を指す識別子の集中は別の形の乖離である。
\textbf{台帳上で標本を組むときは、複製の集中を先に数える必要がある。}
\end{remark}

\begin{remark}[$\Phi$3：対照が識別されない]
\label{rem:phi3-illposed}
$\CCC<0$ は前受けであり、系\ref{cor:growth-needs-capital}によれば $E>0$ を要する。
しかし比較対象である企業部門の $\CCC$ は、
第\ref{part:industry}章の 58 業種すべてで非負である（最小 4.9 日）。
負の $\CCC$ は個別企業に現れるもので、業種の集計では相殺されて消える。

予測と対抗仮説が同じことを述べており、観測しても区別がつかない。
第\ref{sec:predB-illposed}節の可測性仮説と同型の不良設定である。
加えて区分「内」の $\tau$ は秒の単位であり、企業部門の日の単位とは桁が五つ違う。
\end{remark}

\begin{remark}[$\Phi$4：実装できることと現れることは別である]
\label{rem:implementable-not-observed}
族7 は 52 件中 1 件であり、稀であるという量的な部分は当たった。
しかし現れたのは 7-2（クロスサブシディ）であり、7-4 は 0 件であった。

予測は命題\ref{prop:beneficiary}から「7-4 のみが残る」を導いていたが、
命題が述べるのは実装の可否であって出現の頻度ではない。
\textbf{実装できることから現れることを導いたのが誤りである。}

なお唯一の 7-2 は、受益者の選定を署名により手続きの外で行っている。
命題\ref{prop:beneficiary}と整合するが、
予測の文言が外れている以上これを支持の証拠として扱わない。
\end{remark}

\begin{remark}[判定順序の前に一段が要る]
\label{rem:is-it-phi}
表\ref{tab:family-assignment-order}は $\Phi$ を所与としている。
台帳の識別子に当てると、標本 200 件のうち 131 件が $\Phi$ ではなかった。
トークンの台帳（決済の媒体であり $\delta$ と $\pi$ の対を持たない）が 74 件、
対価を伴わない基盤が 27 件、
そして\textbf{決済媒体の名を模した識別子が 30 件}である。

最後のものは定義\ref{def:protocol}を満たすが履行を持たない。
第\ref{sec:identifier}節の $\iota$ の観測不能が、
名称の詐称という形で現れている。
\end{remark}

\section{副次的に測れた量}
\label{sec:protocol-byproducts}

検定の過程で、本稿が他の章で測定を断念していた量が測れた。

\subsection{未想定状態への応答の率。}
命題\ref{prop:unmanned-period}は無人で運転できる期間を $1/\lambda_{\rm novel}$ とした。
$\lambda_{\rm novel}$ は本稿のどこでも測定されていない。

区分「内」では二つの率が測れる。

\begin{center}
\begin{tabularx}{\textwidth}{@{}lXll@{}}
\toprule
量 & 内容 & $\lambda$（毎年） & $1/\lambda$ \\
\midrule
改修の率 $\lambda_{\rm fix}$ & 実装が差し替えられた率。193 万件、うち複製の集中を除く & 0.048 & 20.6 年 \\
停止の率 $\lambda_{\rm stop}$ & 稼働が絶えた率。1,718 件の全数 & 0.030 & 33.8 年 \\
\bottomrule
\end{tabularx}
\captionof{table}{区分「内」で測った応答の率。}
\label{tab:lambda-measured}
\end{center}

\begin{remark}[測っているのは $\lambda_{\rm novel}$ ではない]
\label{rem:lambda-is-response}
表\ref{tab:lambda-measured}の二つは、いずれも\textbf{未想定状態の到来}ではなく
\textbf{到来に対する応答}を数えている。改修は運営者が応答した場合であり、
停止は応答しなかったか、応答できなかった場合である。

到来しても応答を要さなかった場合は、どちらにも現れない。したがって
\[
\lambda_{\rm novel}\ \geq\ \lambda_{\rm fix}+\lambda_{\rm stop}
\]
であり、測れたのは下限である。母集団が異なるため二つは加法的でないが、
\textbf{命題\ref{prop:unmanned-period}の $1/\lambda_{\rm novel}$ は 20 年より短い}という上限が得られる。

改修までの期間の分布は極端に偏る。改修された 20,732 件の中央値は 42 日であるのに対し、
改修されないまま観測を終えた 190 万件の中央値は 4.6 年である。
\textbf{手を入れるものは直ちに入れ、入れないものは何年も入れない。}
平均としての $1/\lambda$ はこの二峰性を覆い隠す。
\end{remark}

\subsection{手続きの継続期間。}
第\ref{ch:platformfee}章は無人継続期間の直接測定を「母集団フレームを欠く」として断念した。
第\ref{sec:protocol-frame}節のとおり、区分「内」ではこの断念の理由が消える。

稼働した 1,718 件の全数について、停止したものが 143 件（8.3\%）、
継続中が 1,575 件である。停止したものの継続期間は中央値 1.62 年、
四分位は 0.64--2.82 年、最大 6.0 年であった。継続中のものは中央値 2.75 年である。
停止したものを事後に集めたのではなく全数から数えているため、
第\ref{sec:survivorship}節の生存条件づけは生じない。

種別ごとの停止率には順序がある。役務 $1/\lambda=11.1$ 年、
高倍率の運用 11.5 年、発行支援 15.3 年、指数 18.1 年、橋渡し 20.1 年、
保険 22.3 年、貸付 22.6 年、利回り 22.7 年である。
\textbf{他の手続きへの依存が薄い種別ほど長い}という向きは読めるが、
予測$\Phi$1 が棄却された以上、この順序を依存の数で説明することはしない。

\subsection{拘束された状態 $B$。}
表\ref{tab:protocol-quantities}は $B$ について「対応する量を本稿は持たない」と記していた。
区分「内」では手続き上に置かれた資産として直接観測でき、
1,718 件の日次系列が得られる。予測$\Phi$2 で $M(t)$ の代理としたのはこれである。

\subsection{料率。}
$B$ が測れると、年間の手数料を $B$ で割った料率が定義できる。
ただし\textbf{分母は族によって異なる}。族3（資産の時間貸し）では $B$ が正しい分母だが、
族6（媒介）では取引額が分母であり $B$ は在庫にすぎない。
混ぜると比が発散する。

族3系に限り、$B$ が 100 万ドル以上のもの 144 件で測った料率は次のとおりである。

\begin{center}
\begin{tabularx}{\textwidth}{@{}Xrrr@{}}
\toprule
種別 & $n$ & 中央値 & 四分位 \\
\midrule
裁定 & 7 & 4.97\% & 2.00--9.31\% \\
実物資産 & 20 & 3.85\% & 2.35--5.84\% \\
担保付債務 & 15 & 3.53\% & 0.99--11.86\% \\
流動化した預託 & 19 & 3.12\% & 2.58--4.21\% \\
貸付 & 30 & 2.97\% & 1.77--4.68\% \\
危険の管理 & 21 & 2.83\% & 2.02--4.47\% \\
利回り & 15 & 2.72\% & 0.76--8.01\% \\
利回りの集約 & 14 & 1.69\% & 0.90--4.19\% \\
\bottomrule
\end{tabularx}
\captionof{table}{族3系の料率（年間手数料 $\div$ 拘束された状態）。}
\label{tab:protocol-fee-ladder}
\end{center}

全体の中央値は 3.07\%、四分位は 1.77--5.84\% である。
第\ref{ch:platformfee}章の料率の梯子と同型の並びだが、
\textbf{比較はできない}。第\ref{ch:platformfee}章の料率は取引額に対する率であり、
ここでの料率は残高に対する率である。次元が違う。

